% R7-06 additive theorem artifact, 2026-09-17 UTC.
% This is a restricted theorem for the I0114 / omega*_9 program.
% It does NOT prove the full striped-theta inequality, omega*_9 universal
% optimality, or Serfaty Conjecture 11.1.

\section*{A rigorous horizontal-axis theorem for the striped index-three blocker}

Let
\[
 e_1=(1,0),\qquad e_2=(1/2,\sqrt3/2),\qquad
 \Lambda=\langle e_1,e_2\rangle,
\]
and let
\[
 M_0=\langle e_1,3e_2\rangle,\qquad
 H=\Lambda\setminus M_0.
\]
For $b>0$ define
\[
 T_b(X,Y)=\sum_{h\in H}\exp\bigl(-b\lvert (X,Y)+h\rvert^2\bigr).
\]
The full near-orbit blocker in the R7-06 reduction is
\[
 T_b(X,Y)\ge T_b(0,0)\qquad (b>0,\ (X,Y)\in\mathbb R^2).
\]
Here we prove this completely on the entire horizontal symmetry axis $Y=0$.

For $c>0$ write the one-dimensional Gaussian comb
\[
 \vartheta_c(t)=\sum_{n\in\mathbb Z}e^{-c(n+t)^2}.
\]
It is positive, even and $1$-periodic.  Put
\[
 U_0=\vartheta_b(0),\qquad U_1=\vartheta_b(1/2),
\]
and
\[
 A=2\vartheta_{27b}(1/3),\qquad
 B=2\vartheta_{27b}(1/6).
\]
The rectangular decomposition already obtained for $H$ gives, at $Y=0$,
\[
 \boxed{T_b(X,0)=A\vartheta_b(X)+B\vartheta_b(X+1/2).}\tag{1}
\]
Also
\[
 T_b(0,0)=A U_0+B U_1.\tag{2}
\]

\subsection*{Lemma 1: a half-shift product inequality}
For every $c>0$ and $t\in\mathbb R$,
\[
 \boxed{\vartheta_c(t)\vartheta_c(t+1/2)
 \ge \vartheta_c(0)\vartheta_c(1/2).}\tag{3}
\]
Moreover, $\vartheta_c(t)\le\vartheta_c(0)$.

\paragraph{Proof.}
Set $q=e^{-\pi^2/c}$.  Poisson summation gives
\[
 \vartheta_c(t)=\sqrt{\frac\pi c}\,\theta_3(\pi t,q),\qquad
 \vartheta_c(t+1/2)=\sqrt{\frac\pi c}\,\theta_4(\pi t,q).
\]
The Jacobi duplication identity is
\[
 \theta_3(z,q)\theta_4(z,q)
 =\theta_4(0,q^2)\theta_4(2z,q^2).\tag{4}
\]
For real $w$, the Jacobi triple product writes
\[
 \theta_4(w,Q)=\prod_{n\ge1}(1-Q^{2n})
 \bigl(1-2Q^{2n-1}\cos(2w)+Q^{4n-2}\bigr).
\]
Every factor depending on $w$ is minimized when $\cos(2w)=1$, hence
\[
 \theta_4(w,Q)\ge\theta_4(0,Q).
\]
Applying this to (4) and then evaluating (4) at $z=0$ yields (3).
Finally, the Poisson Fourier series for $\vartheta_c$ has positive cosine
coefficients, so its maximum is attained at an integer; thus
$\vartheta_c(t)\le\vartheta_c(0)$. \qed

\subsection*{Lemma 2: the conductor-$27$ coefficient ratio}
For every $b>0$,
\[
 \boxed{
 \frac{\vartheta_{27b}(1/6)}{\vartheta_{27b}(1/3)}
 >\frac{\vartheta_b(0)}{\vartheta_b(1/2)}.}\tag{5}
\]

\paragraph{Proof.}
Let
\[
 r=e^{-\pi^2/(27b)}\in(0,1).
\]
Poisson summation converts the two sides of (5) to
\[
 L(r):=\frac{\theta_3(\pi/6,r)}{\theta_3(\pi/3,r)},\qquad
 R(r):=\frac{\theta_3(0,r^{27})}{\theta_4(0,r^{27})}.
\]
For $0<x<1$ define
\[
 h(x)=\log\frac{1+x}{1-x}.
\]
Then $h$ is positive and strictly increasing.  In the triple product for
$\theta_3$, the ratio between the factors at $\pi/6$ and $\pi/3$, for an odd
exponent $m=2n-1$, is
\[
 \frac{1+r^m+r^{2m}}{1-r^m+r^{2m}}
 =\frac{1+r^m}{1-r^m}\frac{1-r^{3m}}{1+r^{3m}}.
\]
Therefore, with absolutely convergent sums,
\[
 \begin{aligned}
 \log L(r)
 &=\sum_{m\ge1\atop m\ {\rm odd}}
   \bigl(h(r^m)-h(r^{3m})\bigr)\\
 &=\sum_{k\ge1\atop k\ {\rm odd},\ 3\nmid k}h(r^k).\tag{6}
 \end{aligned}
\]
Similarly the triple products at $0$ give
\[
 \log R(r)=2\sum_{m\ge1\atop m\ {\rm odd}}h(r^{27m}).\tag{7}
\]
For every positive odd $m$, select the two indices
\[
 k_1(m)=27m-2,\qquad k_2(m)=27m-4.
\]
They are positive odd integers, neither is divisible by $3$, and all selected
indices are pairwise distinct (their residues modulo $27$ are $25$ and $23$).
Moreover $k_i(m)<27m$, hence, because $0<r<1$ and $h$ is strictly increasing,
\[
 h(r^{k_1(m)})+h(r^{k_2(m)})
 >2h(r^{27m}).\tag{8}
\]
The left side of (8), summed over odd $m$, is only a subseries of the positive
sum in (6).  Comparing with (7) proves $\log L(r)>\log R(r)$, and hence (5).
\qed

\subsection*{Theorem: global minimum on the horizontal axis}
For every $b>0$ and $X\in\mathbb R$,
\[
 \boxed{T_b(X,0)\ge T_b(0,0).}\tag{9}
\]
Equality holds precisely when $X\in\mathbb Z$.

\paragraph{Proof.}
Set
\[
 u=\vartheta_b(X),\qquad v=\vartheta_b(X+1/2),\qquad p=U_0U_1.
\]
By Lemma 1,
\[
 uv\ge p,\qquad u\le U_0.\tag{10}
\]
By Lemma 2,
\[
 \frac BA>\frac{U_0}{U_1},\qquad\text{equivalently}\qquad
 \frac{Bp}{U_0^2}>A.\tag{11}
\]
Using $v\ge p/u$ in (1),
\[
 T_b(X,0)\ge Au+\frac{Bp}{u}=:g(u).\tag{12}
\]
For every $0<u\le U_0$,
\[
 g'(u)=A-\frac{Bp}{u^2}
 \le A-\frac{Bp}{U_0^2}<0.\tag{13}
\]
Thus $g$ is strictly decreasing on the entire permitted interval for $u$, and
because $u\le U_0$,
\[
 T_b(X,0)\ge g(u)\ge g(U_0)
 =AU_0+BU_1=T_b(0,0).
\]
If equality holds, (13) forces $u=U_0$, hence $X\in\mathbb Z$; conversely
$T_b$ is $M_0$-periodic and $e_1\in M_0$, so every integral $X$ gives equality.
\qed

\subsection*{Interpretation for the near rigid-pair orbit}
In the R7-06 rigid-pair reduction, the original Gaussian scale and common
translation are related by $b=a/9$ and $y=3t$.  The theorem therefore proves
that the near pair $\{0,e_1/3\}$ cannot lower the Gaussian energy under any
common translation parallel to $e_1$.  Integrating the Gaussian inequality
against a positive Bernstein measure extends this restricted theorem to every
source-admissible completely monotone function of squared distance for which
the periodic energy is finite.

This is a genuine two-defect theorem but only a one-dimensional slice of the
full striped inequality.  It does not move the frozen source metric $m_{\rm sq}$
from $2$ and does not prove global universal optimality of $\omega^*_9$.
